\documentclass[14pt,usepdftitle=false,aspectratio=169]{beamer}
\usepackage{preambule}
\setbeamercolor{structure}{fg=black}
\usepackage{polynomes,arithmetique,arrows,al,matrices}\DeclareMathOperator{\olddisc}{disc}\newcommand{\disc}[1]{\olddisc\l#1\r}
\hypersetup{pdftitle={Théorie algébrique des nombres -- Formes quadratiques binaires}}
\title{Théorie algébrique des nombres\\\emph{Formes quadratiques binaires}}
\author{}
\date{}
\begin{document}
\begin{frame}
    \titlepage
\end{frame}
\slideq{CNS sur $\disc q$ pour que $q$ représente primitivement $n$}{1/22}
\slider{$\cgr{\disc q}0{4n}$}{1/22}
\slideq{Forme principale de détermiant $D$}{2/22}
\slider{$q_D=\l1,0,\frac{-D}4\r$ si $\cgr D04$\linebreak$q_D=\l1,1,\frac{1-D}4\r$ si $\cgr D14$}{2/22}
\slideq{CNS sur $\disc q$ pour que $q$ représente $1$}{3/22}
\slider{$q\sim^+q_{\disc q}$ la forme principale de détermiant $\disc q$}{3/22}
\slideq{CNS pour que $q$ représente $0$}{4/22}
\slider{$D$ est un carré}{4/22}
\slideq{Forme quadratique binaire primitive}{5/22}
\slider{$q=\l a,b,c\r$ avec $a\wedge b\wedge c=1$}{5/22}
\slideq{$q\sim q'$}{6/22}
\slider{$\exists P\in\matgl2\bbZ,q'=q\circ P$}{6/22}
\slideq{Forme quadratique binaire}{7/22}
\slider{$q\in\pol Z{X,Y}$ homogène de degré $2$\linebreak$\l a,b,c\r=aX^2+bXY+cY^2$}{7/22}
\slideq{$\disc q$\linebreak$q=\l a,b,c\r$}{8/22}
\slider{$D=b^2-4ac$}{8/22}
\slideq{CNS sur la classe de $q$ pour qu'elle représente primitivement $n\in\bbZ$}{9/22}
\slider{Il existe $b$ tel que $-\left|n\right|<b\leqslant\left|n\right|$ et $c\in\bbZ$ tels que $q\sim^+\l n,b,c\r$}{9/22}
\slideq{$q\sim^+ q'$}{10/22}
\slider{$\exists P\in\sl[2]\bbZ,q'=q\circ P$}{10/22}
\slideq{$M_q$\linebreak$q=\l a,b,c\r$}{11/22}
\slider{$\crampedtmatrix({2a\&b\\b\&2c\\})$\linebreak$q\l x,y\r=\frac12\crampedtmatrix({x\\y\\})^\top M_q\crampedtmatrix({x\\y\\})$}{11/22}
\slideq{$q$ représente (primitivement) $n$}{12/22}
\slider{Il existe $\l x,y\r\in\bbZ^2\setminus\set{\l0,0\r}$ ($x\wedge y=1$) tel que $q\l x,y\r=n$}{12/22}
\slideq{CNS pour que $D\leqslant0$}{13/22}
\slider{$q$ est de signe constant}{13/22}
\slideq{Lien entre $\operatorname{Cl}\l\calO_K\r$ et $\operatorname{Cl}\l D\r$ pour $K=\fr Q{\sqrt d}$ avec $d<0$}{14/22}
\slider{$I\mapsto\l\l x,y\r\mapsto\frac{\operatorname{nrm}_{K/\bbQ}\l xu+yv\r}{N\l I\r}\r$ pour $\l u,v\r$ une base directe de $I$\linebreak $q\mapsto\bbZ a\oplus\bbZ  a\tau$ pour $\tau=\frac{b\pm\sqrt D}{2a}$ choisi tel que $\l a,a\tau\r$ soit directe\linebreak Ces deux applications sont deux bijections réciproques}{14/22}
\slideq{Théorème de Gauss}{15/22}
\slider{Si $D<0$, toute forme quadratique binaire positive est proprement équivalente à une unique forme réduite}{15/22}
\slideq{Théorème de Lagrange}{16/22}
\slider{Si $D\neq0$, il n'existe qu'un nombre fini de classes de $\calQ\l D\r$ par l'action de $\sl[2]\bbZ$}{16/22}
\slideq{Base directe de $I\leqslant\calO_K$, $K=\fr Q{\sqrt d}$}{17/22}
\slider{$\l u,v\r=P\l1,\alpha\r$ avec $P\in\matgl[+]2\bbQ$ et $\alpha=\sqrt d$ si $\cgr d24$ ou $\cgr d34$ et $\alpha=\frac{1+\sqrt d}{2}$ si $\cgr d14$}{17/22}
\slideq{$q\in\calQ$ est positive}{18/22}
\slider{$q\l x,y\r\geqslant0$ pour tout $\l x,y\r\in\bbZ^2$}{18/22}
\slideq{Action $\matgl2\bbZ\acts\calQ$}{19/22}
\slider{$q\circ P\l x,y\r=q(P(x,y))$\linebreak$M_{q\circ P}=P^\top M_qP$\linebreak L'action préserve le discriminant, les formes primitives, les entiers (primitivement) représentés et le signe (si $D\leqslant0$)}{19/22}
\slideq{Condition pour que $p$ soit représentée par $q$ de discriminant $D$ et propriétés de ces formes quadratique}{20/22}
\slider{Si $D$ est un carré modulo $4p$, il existe une forme quadratique binaire représentant $p$\linebreak Elle est unique à équivalence près par l'acion de $\matgl2\bbZ$}{20/22}
\slideq{CNS pour que $D$ soit le discriminant d'une forme quadratique binaire}{21/22}
\slider{$\cgr D04$ ou $\cgr D14$}{21/22}
\slideq{La forme positive $q=\l a,b,c\r$ est réduite}{22/22}
\slider{$-a<b\leqslant a\leqslant c$ avec $b\geqslant0$ si $a=c$}{22/22}
\end{document}